% ================================================================
\section{Conclusioni}\label{sec:conclusioni}

L'obiezione va formulata nella sua forma più forte: tutto ciò che la
\S\ref{sec:architettura} descrive potrebbe girare in tempo reale su un
laptop, e un granulatore differito ha oggi l'aria dell'anacronismo. La
scelta richiede perciò un argomento, non una giustificazione tecnica ---
e l'argomento ha un precedente esplicito, formulato quando la posizione
era già controcorrente. Per Risset la composizione \emph{«is not --- or
should not be --- a real-time process»}: la notazione riferisce la
musica a una rappresentazione fuori dal tempo, e proprio l'operare non
in tempo reale libera dalla \emph{«arrow of time and its tyranny»},
dai riflessi e dalla fretta; \emph{«writing music implies prediction and
elaboration»}~\cite{Risset1999}. Non è nostalgia del nastro: è la tesi
che un certo tipo di lavoro compositivo --- predizione, elaborazione,
ritorno sul già scritto --- abbia nel differito la propria
configurazione naturale.








